\clearpage
\fancyhead[L]{\small GAUSS AN VON LINDENAU.}
\SourceLeaf{310}\begin{center}
GAUSS AN VON LINDENAU.
\end{center}

\hrule
\begin{center}
\fontsize{10}{12}\selectfont Monatliche Correspondenz. Band XVIII. S. 173..175. 1808 August.
\end{center}
\hrule

\begin{flushright}
Göttingen, den 6. August 1808.
\end{flushright}

\par\indent Seit meinem letzten Briefe habe ich die Beobachtungen der neuen Planeten, so oft die Umstände
es erlaubten, fortgesetzt, und ich mache mir das Vergnügen Ihnen hier mitzutheilen, was ich bisher
erhalten habe. \emph{Juno}, \emph{Vesta} und \emph{Pallas} konnten bisher nur mit dem Kreismicrometer beobachtet werden.
Die Beobachtungen der Vesta habe ich aber seit dem 13. Juli nicht fortgesetzt, da sie von jetzt an schon
im Meridian wird beobachtet werden können. Von der \emph{Ceres} ist Hrn. Prof. \textsc{Harding} auch eine Beobachtung am Quadranten gelungen, die ich Ihnen künftig mittheilen werde [handschriftliche Bemerkung:
Juli 25. $13^{\text{u}}\,14^{\text{m}}\,29^{\text{s}}$ \quad $322^\circ\,8'\,57''6$ \quad $27^\circ\,52'\,11''5$]

\par\indent Die \emph{Pallas} steht jetzt leider die ganze Nacht für unser Local zu hoch; ich habe nur einige Beobachtungen erhalten können, wobei aber die Declinationen zweifelhaft sind.

\par\indent Die Beobachtung der \emph{Vesta} vom 2. Juli und die der \emph{Pallas} vom 15. Juli sind von Hrn. \textsc{Harding},
die übrigen alle von mir. Der Fehler der Ephemeride bei der \emph{Pallas} ist $-4'$ in gerader Aufsteigung,
die Declination auf die $1'$ richtig; die Differenz bei der \emph{Juno} ist jetzt auf $13'$ in gerader Aufsteigung
angewachsen. Die vier letzten Beobachtungen der \emph{Juno} halte ich alle für sehr gut, in so fern die verglichenen Sterne in der Hist. Cél. gut beobachtet sind.

\Needspace{13\baselineskip}
\begin{center}
\emph{Beobachtungen der Pallas.}
\end{center}


\begin{center}\fontsize{10}{12}\selectfont\setlength{\tabcolsep}{3pt}\renewcommand{\arraystretch}{1.15}\begin{tabular}{lrrr}
1808 & \shortstack{Mittlere Zeit\\in Göttingen} & \shortstack{Scheinbare\\AR.} & \shortstack{Scheinbare\\Declination}\\\hline
Juli 13 & $12^{\text{u}}\;3^{\text{m}}\;17^{\text{s}}$ & $300^\circ\;21'\;49''$ & $18^\circ\;41'\;44''$ N.\\
15 & $12\;1\;28$ & $299\;57\;3$ & $18\;32\;40$\\
25 & $10\;9\;44$ & $297\;57\;37$ & $17\;40\;10$\\
\end{tabular}\end{center}


\Needspace{13\baselineskip}
\begin{center}
\emph{Beobachtungen der Juno.}
\end{center}


\begin{center}\fontsize{10}{12}\selectfont\setlength{\tabcolsep}{3pt}\renewcommand{\arraystretch}{1.15}\begin{tabular}{lrrr}
1808 & \shortstack{Mittlere Zeit\\in Göttingen} & \shortstack{Scheinbare\\AR.} & \shortstack{Scheinbare\\Declination}\\\hline
Juni 20 & $11^{\text{u}}\;49^{\text{m}}\;0^{\text{s}}$ & $315^\circ\;29'\;34''$ & $2^\circ\;16'\;23''$ S.\\
22 & $12\;0\;45$ & $315\;23\;1$ & $2\;14\;29$\\
Juli 6 & $12\;42\;23$ & $314\;1\;0$ & $2\;19\;7$\\
30 & $10\;51\;17$ & $309\;38\;9$ & $3\;57\;13$\\
31 & $11\;40\;57$ & $309\;24\;53$ & $4\;3\;48$\\
Aug. 4 & $10\;42\;53$ & $308\;33\;44$ & $4\;30\;32$\\
5 & $10\;53\;46$ & $308\;20\;39$ & $4\;37\;36$\\
\end{tabular}\end{center}


\par\indent [Im Original sind noch Beobachtungen der \emph{Vesta} von Juni 22. Juli 1. 2. 13., beigefügt, die sich
aber auch in der weiter unten S. 309 abgedruckten Mittheilung der Gött. gel. Anz. 1808 Nov. 10 finden.]

\par\indent Ich habe bereits levi calamo aus meinen Beobachtungen die Opposition der \emph{Juno} berechnet und
die Elemente vorläufig verbessert. Eine nur kleine Änderung, die vornehmlich nur in einer kleinen
Verminderung der Länge der Sonnennähe und einer kleinen Vergrösserung der mittlern Bewegung besteht, reicht hin, alle vier bisher beobachteten Oppositionen in einer reinen Ellipse sehr genau darzustellen. Nächstens hievon mehr.

\par\indent Unsere Sternwarte hat von dem Könige mehrere astronomische Instrumente zum Geschenk erhalten, wovon ein Gregorianisches Telescop von \textsc{Chevalier} und eine Seeuhr von \textsc{Louis Berthoud} bereits an\SourceLeaf{311}gekommen sind. Letztere hat, so viel sich aus den bisherigen Erfahrungen schliessen lässt, einen ziemlich gleichförmigen Gang. Freilich ist derselbe bisher nur bei beständiger Ruhe beobachtet worden.
Nächstens werden wir den Versuch einer Längenbestimmung damit machen. [Diese Mittheilung über
die Instrumente ist auch in den Gött. gel. Anzeigen Stück 127. Seite 1265. 1808 August 8. abgedruckt
mit dem Zusatze „Die übrigen Instrumente welche bisher in Paris befindlich waren, dürfen wir gleichfalls bald erwarten.“ \textsc{Schering}.]

\Needspace{15\baselineskip}
\fancyhead[L]{\small JUNO ELEMENTE VIII.}
\begin{center}
JUNO ELEMENTE VIII.
\end{center}

\hrule
\begin{center}
\fontsize{10}{12}\selectfont Göttingische gelehrte Anzeigen. Stück 136. Seite 1353..1355. 1808 August 25.
\end{center}
\hrule

\par\indent Die ersten, in diesem Jahre auf der hiesigen Sternwarte gemachten Beobachtungen der \emph{Juno} sind
bereits im 107. Stück dieser Blätter [S. 303 d. B.] angezeigt: seitdem sind die Beobachtungen fortgesetzt,
so oft die Umstände dazu günstig waren, freilich nur mit dem Kreismicrometer, da der kleine Planet
in dem lichtschwachen Fernrohre unsers Mauerquadranten beständig unsichtbar blieb. Man hat diesen
Mangel durch zahlreiche und möglichst sorgfältige Vergleichungen zu ersetzen gesucht, und besonders
die vier Beobachtungen in der Nähe der Opposition dürfen für so genau gehalten werden, als es nur
immer die angewandte Methode verstattet. Da die Beobachtung vom 22. Juni noch eine kleine Correction erlitten hat, so stellen wir hier die ersten beiden Bestimmungen noch einmal mit allen bisher
gemachten späteren zusammen: [stehen oben Seite 305 in dem Briefe an \textsc{von Lindenau}. 1808 Aug. 6.]

\par\indent Der Fehler der Ephemeride ist hienach bei den letzten Beobachtungen auf 13 Minuten in gerader Aufsteigung angewachsen; der Fehler der Declination ist ziemlich unverändert $1\frac{1}{2}$ Minute. Auswärtige Beobachtungen sind bisher noch nicht bekannt geworden.

\par\indent Hr. Prof. \textsc{Gauss} hat die vier letzten, vorzüglich gut ausgefallenen Beobachtungen zur Bestimmung der Opposition benutzt, und folgendes Resultat gefunden:

\begin{center}
1808. August 2. $9^{\text{u}}\,30^{\text{m}}\,43^{\text{s}}$ mittlere Zeit in Göttingen\\
wahre Länge .\;.\;.\;.\;.\;.\;.\;.\;.\; $310^\circ\,16'\,31''5$\\
wahre geocentrische Breite .\;.\;.\;.\; $13\;53\;56,2$ nordl.
\end{center}

\par\indent Die Verbindung dieser Opposition, der \emph{vierten} bisher beobachteten, mit denen von 1804, 1806 und
1807 hat hienächst zur Bestimmung folgender neuen Elemente (VIII) gedient, wodurch die sämmtlichen
bisherigen Beobachtungen noch sehr gut dargestellt werden.


\begin{center}\fontsize{10}{12}\selectfont\setlength{\tabcolsep}{3pt}\renewcommand{\arraystretch}{1.15}\begin{tabular}{lr}
\multicolumn{2}{c}{Epoche der mittlern Länge}\\
\multicolumn{2}{c}{für den Meridian von Göttingen:}\\\hline
1804 .\;.\;.\;.\;.\;.\;.\;.\; & $320^\circ\;1'\;20''1$\\
1805 .\;.\;.\;.\;.\;.\;.\;.\; & $42\;35\;8.4$\\
1806 .\;.\;.\;.\;.\;.\;.\;.\; & $125\;8\;56.7$\\
1807 .\;.\;.\;.\;.\;.\;.\;.\; & $207\;42\;45.0$\\
1808 .\;.\;.\;.\;.\;.\;.\;.\; & $290\;30\;7.6$\\
1809 .\;.\;.\;.\;.\;.\;.\;.\; & $13\;3\;55.8$\\
1810 .\;.\;.\;.\;.\;.\;.\;.\; & $95\;37\;44.1$\\
\end{tabular}\end{center}

\begin{center}\fontsize{9}{11}\selectfont\setlength{\tabcolsep}{3pt}\renewcommand{\arraystretch}{1.15}\begin{tabular}{lr}
Tägliche mittlere tropische Bewegung .\;.\; & $814''324$\\
Tropische Umlaufszeit .\;.\; & 1591 Tage 12 Stunden\\
Sonnennähe 1805 .\;.\;.\;.\;.\;.\;.\; & $53^\circ\;10'\;53''9$\\
Aufsteigender Knoten 1805 .\;.\;.\;.\; & $171\;4\;11{,}3$\\
\multicolumn{2}{l}{\quad Beide siderisch ruhend vorausgesetzt.}\\
Neigung der Bahn .\;.\;.\;.\;.\;.\;.\;.\; & $13\;4\;11{,}0$\\
Excentricität .\;.\;.\;.\;.\;.\;.\;.\;.\;.\;.\; & $0{,}2554521$\\
Logarithm der halben grossen Axe .\;.\;.\; & $0{,}4261883$\\
\end{tabular}\end{center}


\SourceLeaf{312}\par\indent Die nächste Opposition, wo die Juno wieder ansehnlich heller sein wird, fällt nach diesen Elementen 1810 Januar 30; in $130^\circ\,2'$ Länge und $14^\circ\,51'$ südlicher Breite, am Kopf der Wasserschlange.

\Needspace{16\baselineskip}
\fancyhead[L]{\small GAUSS AN VON LINDENAU.}
\begin{center}
GAUSS AN VON LINDENAU.
\end{center}

\hrule
\begin{center}
\fontsize{10}{12}\selectfont Monatliche Correspondenz. Band XVIII. S. 269..273. 1808 September.
\end{center}
\hrule

\begin{flushright}
Göttingen, am 30. August 1808.
\end{flushright}

\par\indent Auf das verbindlichste danke ich Ihnen für die gütige Mittheilung der Nachrichten von den beiden letzten Cometen und Ihrer Beobachtungen der neuen Planeten. Es ist Schade, dass erstere nicht
früher bekannt geworden sind, jetzt ist natürlich an die Aufsuchung nicht mehr zu denken. Die letzteren Resultate der hiesigen Juno-Beobachtungen erlauben Sie mir, Ihnen auf beiliegendem Blatte unserer gel. Anz. vorlegen zu dürfen. Die Beobachtungen waren folgende:


\begin{center}\fontsize{10}{12}\selectfont\setlength{\tabcolsep}{3pt}\renewcommand{\arraystretch}{1.15}\begin{tabular}{lrrr}
1808 & \shortstack{Mittlere Zeit\\in Göttingen} & \shortstack{Scheinb. gerade\\Aufst. der Juno} & \shortstack{Südl. Abweich.\\der Juno}\\\hline
Juni 20 & $11^{\text{u}}\;49^{\text{m}}\;0^{\text{s}}$ & $315^\circ\;29'\;34''$ & $2^\circ\;16'\;23''$\\
22 & $12\;0\;45$ & $315\;23\;1$ & $2\;14\;29$\\
Juli 6 & $12\;42\;23$ & $314\;1\;0$ & $2\;19\;7$\\
30 & $10\;51\;17$ & $309\;38\;0$ & $3\;57\;13$\\
31 & $11\;40\;57$ & $309\;24\;53$ & $4\;3\;48$\\
Aug. 4 & $10\;42\;53$ & $308\;33\;44$ & $4\;30\;32$\\
5 & $10\;53\;46$ & $308\;20\;39$ & $4\;37\;36$\\
\end{tabular}\end{center}


\Needspace{11\baselineskip}
\par\indent Der Fehler der Ephemeride ist hienach bei den letzten Beobachtungen auf $13'$ in gerader Aufsteigung angewachsen; der Fehler der Declination ist ziemlich unverändert $1'5$. Auswärtige Beobachtungen sind bisher noch nicht bekannt geworden. Die vier letzten Beobachtungen geben für die Opposition folgendes Resultat:

\begin{center}
1808 August 2. $9^{\text{u}}\,30^{\text{m}}\,43^{\text{s}}$ mittlere Zeit in Göttingen\\
Wahre Länge .\;.\;.\;.\;.\;.\;.\;.\;.\; $310^\circ\,16'\,31''5$\\
Wahre geocentrische Breite .\;.\;.\;.\; $13\;53\;56,5$ nordl. \end{center}

\par\indent Die Verbindung dieser Opposition, der vierten bisher beobachteten, mit denen von 1804, 1806 und
1807, hat hienächst zur Bestimmung folgender neuen Elemente (VIII) gedient, wodurch die sämmtlichen
bisherigen Beobachtungen noch sehr gut dargestellt werden, [sie sind oben Seite 306 aus den Gött. gel.
Anzeigen 1808 Aug. 25 abgedruckt.]

\par\indent Die Vesta haben wir schon achtmal seit dem 12. August im Meridian beobachtet; die Beobachtungen sind aber noch nicht alle berechnet.

\par\indent [Den fernern Theil des Briefes enthaltend die Ankündigung des Aufsatzes über die Bestimmung
der Polhöhe habe ich oben nach jenem Aufsatze Seite 146 dieses Bandes wiedergegeben. \textsc{Schering}.]
